\documentclass[11pt]{article}

\usepackage[a4paper,margin=2.5cm]{geometry}
\usepackage{amsmath,amssymb,amsthm,mathtools}
\usepackage{enumitem}
\usepackage{hyperref}

\title{Radial rigidity for $\Delta u=f(|x|)\mathbf 1_U-c\,\delta_0$}
\author{}
\date{}

\newtheorem{theorem}{Theorem}[section]
\newtheorem{proposition}[theorem]{Proposition}
\newtheorem{lemma}[theorem]{Lemma}
\newtheorem{corollary}[theorem]{Corollary}
\newtheorem{remark}[theorem]{Remark}

\newcommand{\C}{\mathbb C}
\newcommand{\R}{\mathbb R}
\newcommand{\1}{\mathbf 1}

\begin{document}
\maketitle

\section{Statement}

Let $U\subset\R^2\simeq\C$ be a bounded open set with $C^1$ boundary, and let
\[
f:(0,\infty)\to\R
\]
be continuous with $f(r)>0$ for every $r>0$. Let $c\in\R$. Consider the distributional
equation
\begin{equation}\label{eq:main}
\Delta u=f(|x|)\,\1_U-c\,\delta_0\qquad\text{in }\mathcal D'(\R^2).
\end{equation}

\begin{theorem}\label{thm:main}
Assume $0\notin\partial U$. Equation \eqref{eq:main} admits a compactly supported
distributional solution $u$ with $u\in C^1(\R^2\setminus\{0\})$ if and only if
\begin{equation}\label{eq:c-eq}
c=\int_U f(|y|)\,dA(y)
\end{equation}
and $U$ is rotationally invariant about the origin.
\end{theorem}

\begin{remark}
A bounded rotationally invariant open subset of $\R^2$ is a disjoint union of (concentric)
open annuli about the origin, together with at most one open disk centered at the origin.
If one additionally assumes that $U$ is connected, rotational invariance forces $U$ to be a
disk or a single annulus centered at the origin.
\end{remark}

\section{The easy direction: radial $U$ admits a compactly supported solution}

\begin{proposition}\label{prop:converse}
Let $U=\bigsqcup_{j=1}^N A_j$ be a disjoint union of (concentric) open annuli or disks
$A_j=\{r_j^-<|x|<r_j^+\}$ (with the convention $r_j^-\ge 0$; $r_j^-=0$ gives a disk) centered
at the origin, ordered so that $r_1^+<r_2^-\le r_2^+<\cdots$. Assume $f$ is continuous on
$(0,\infty)$ and that \eqref{eq:c-eq} holds. Then \eqref{eq:main} admits a compactly
supported radial solution $u\in C^1(\R^2\setminus\{0\})\cap C^2(\R^2\setminus(\partial U\cup\{0\}))$.
\end{proposition}

\begin{proof}
We look for $u=u(r)$ depending only on $r=|x|$. Set
\[
F(r):=\int_0^r s\,f(s)\,\1_U(s)\,ds,
\]
so that $F$ is continuous, nondecreasing, and $F(r)=\frac{1}{2\pi}c$ for all $r$ larger than
$r_N^+$ by \eqref{eq:c-eq}.

Define, for $r>0$,
\[
g(r):=-\frac{c}{2\pi}+2\pi\int_0^r sf(s)\1_U(s)\,ds\cdot\frac1{2\pi}
      =-\frac{c}{2\pi}+F(r),
\]
and set
\[
u'(r):=\frac{g(r)}{r},\qquad u(r):=-\int_r^{\infty}\frac{g(s)}{s}\,ds.
\]
The integral converges since $g(s)=0$ for $s\ge r_N^+$ by \eqref{eq:c-eq}. Moreover
$u(r)=0$ for $r\ge r_N^+$, so $u$ is compactly supported.

For $r>0$ we compute in the classical sense
\[
\Delta u(r)=\frac1r\bigl(ru'(r)\bigr)'=\frac{g'(r)}{r}=f(r)\1_U(r)
\qquad(r\in(0,\infty)\setminus\{r_j^\pm\}).
\]
Across each $r=r_j^\pm$, $g$ is continuous, hence $u'$ is continuous; so
$u\in C^1((0,\infty))$.

Near $r=0$, $g(r)\to -c/(2\pi)$, so
\[
u(r)=-\frac{c}{2\pi}\log r+O(1)\qquad(r\to 0^+).
\]
Standard computation shows that the distributional Laplacian of $u$ on $\R^2$ is
\[
\Delta u=f(|x|)\1_U-c\,\delta_0,
\]
since the radial profile $r\mapsto -\frac{1}{2\pi}\log r$ has distributional Laplacian
$\delta_0$. This proves the proposition.
\end{proof}

\section{The hard direction: a compactly supported solution forces $U$ to be radial}

The proof follows the scheme suggested in \texttt{connected-annulus-rigidity.tex}, but we
do \emph{not} first show that the outer boundary is a circle. The crucial observation is
that the vanishing of $u$ in the unbounded exterior $E_\infty$ already gives enough Cauchy
data on $\Gamma_{\mathrm{out}}$ to kill the angular derivative of $u$ across the whole
connected component, irrespective of the shape of $\Gamma_{\mathrm{out}}$.

Throughout this section assume:
\begin{enumerate}[label=(H\arabic*)]
\item $U\subset\R^2$ is bounded, open, with $C^1$ boundary, and $0\notin\partial U$;
\item $f:(0,\infty)\to\R$ is continuous with $f(r)>0$ for all $r>0$;
\item $u\in C^1(\R^2\setminus\{0\})$ is compactly supported and
satisfies \eqref{eq:main} distributionally.
\end{enumerate}

\subsection{Preliminary reductions}

Denote by $E_\infty$ the unbounded connected component of $\R^2\setminus\overline U$, and by
$\Gamma_{\mathrm{out}}:=\partial E_\infty$ the outer boundary of $U$. Since $0\notin\partial U$,
$\Gamma_{\mathrm{out}}$ is a compact subset of $\R^2\setminus\{0\}$.

\begin{lemma}\label{lem:c-eq}
Under hypotheses (H1)--(H3), one has
\[
c=\int_U f(|y|)\,dA(y).
\]
\end{lemma}

\begin{proof}
Test \eqref{eq:main} with a cutoff $\chi\in C_c^\infty(\R^2)$ equal to $1$ on a neighborhood
of $\mathrm{supp}(u)\cup\{0\}$. Then $\int u\Delta\chi=0$ since $\chi$ is constant on
$\mathrm{supp}(u)$, and the right-hand side of \eqref{eq:main} paired with $\chi$ gives
$\int_U f-c$.
\end{proof}

\begin{lemma}\label{lem:vanish-Einfty}
Under (H1)--(H3), assume furthermore $0\notin E_\infty$. Then
\[
u\equiv 0\qquad\text{on }E_\infty.
\]
\end{lemma}

\begin{proof}
On $E_\infty$, the right-hand side of \eqref{eq:main} vanishes: $\1_U=0$ there, and
$\delta_0$ is not supported in $E_\infty$ by hypothesis. Hence $\Delta u=0$ in $E_\infty$,
so $u$ is harmonic there (Weyl's lemma). Since $u$ is compactly supported, $u$ vanishes
on the complement of a large ball, hence on a nonempty open subset of $E_\infty$. As
$E_\infty$ is connected and $u$ is real-analytic there, $u\equiv 0$ on $E_\infty$.
\end{proof}

\begin{remark}
If $0\in E_\infty$, then on $E_\infty$ we have $\Delta u=-c\,\delta_0$, so
$u(x)=-\frac{c}{2\pi}\log|x|+h(x)$ with $h$ harmonic on $E_\infty$. But $u$ is compactly
supported, which forces $c=0$ (by Lemma~\ref{lem:c-eq}) and then $h\equiv 0$; the argument
below then applies verbatim. We henceforth tacitly assume $0\notin E_\infty$, which is the
only nontrivial case.
\end{remark}

\subsection{The angular derivative is harmonic}

Let $L:=x\partial_y-y\partial_x=\partial_\theta$ denote the generator of rotations, and set
\[
w:=Lu=\partial_\theta u\in C^0(\R^2\setminus\{0\}).
\]
Note that $w$ need not be globally $C^1$ because $L$ is only a first-order operator and $u$
is a priori only $C^1$. Nevertheless, $w$ is defined pointwise away from $0$.

\begin{lemma}\label{lem:w-harmonic}
Under (H1)--(H3), $w$ is harmonic in the open set
$\Omega:=\R^2\setminus(\partial U\cup\{0\})$.
In particular $w$ is smooth in $\Omega$.
\end{lemma}

\begin{proof}
The operator $L$ commutes with $\Delta$ and with multiplication by any radial function, and
$L$ annihilates both radial functions and $\delta_0$. Therefore, in the interior of $U$
(where $\1_U$ is locally constant $=1$) and in the interior of $\R^2\setminus\overline U$
(where $\1_U=0$), away from the origin, we have
\[
\Delta w=\Delta Lu=L\Delta u=L\bigl(f(|x|)\1_U-c\,\delta_0\bigr)=0.
\]
Weyl's lemma upgrades the distributional identity to $w\in C^\infty(\Omega)$ harmonic.
\end{proof}

\subsection{Cauchy data of $w$ on the outer boundary}

\begin{lemma}\label{lem:cauchy-data}
Under (H1)--(H3), and assuming $0\notin E_\infty$:
\begin{enumerate}[label=(\alph*)]
\item $u\equiv 0$ and $\nabla u\equiv 0$ on $\Gamma_{\mathrm{out}}$;
\item the extension $\tilde w$ of $w$ by $0$ to a two-sided $C^1$ neighborhood
$\mathcal N$ of $\Gamma_{\mathrm{out}}$ (inside $U$ on one side and inside $E_\infty$ on the
other) is continuous, with $\tilde w=0$ on $\mathcal N\cap(E_\infty\cup\Gamma_{\mathrm{out}})$.
\end{enumerate}
\end{lemma}

\begin{proof}
(a) By Lemma~\ref{lem:vanish-Einfty}, $u\equiv 0$ in $E_\infty$, so the limiting values of
$u$ and $\nabla u$ from the exterior side of $\Gamma_{\mathrm{out}}$ vanish. Since
$u\in C^1(\R^2\setminus\{0\})$ and $\Gamma_{\mathrm{out}}\subset\R^2\setminus\{0\}$, the
interior-side limits agree with the exterior-side ones, giving
$u=0$ and $\nabla u=0$ on $\Gamma_{\mathrm{out}}$.

(b) Because $\nabla u=0$ on $\Gamma_{\mathrm{out}}$, the angular derivative
$w=x\partial_y u-y\partial_x u$ vanishes on $\Gamma_{\mathrm{out}}$. On the other side,
$w\equiv 0$ in $E_\infty$ since $u\equiv 0$ there. Thus extending by $0$ across
$\Gamma_{\mathrm{out}}$ yields a continuous function $\tilde w$ on a neighborhood
$\mathcal N$ of $\Gamma_{\mathrm{out}}$.
\end{proof}

\subsection{The key step: $w$ vanishes near $\Gamma_{\mathrm{out}}$}

\begin{lemma}\label{lem:extension-harmonic}
Under the standing hypotheses, the extension $\tilde w$ from Lemma~\ref{lem:cauchy-data}(b)
is harmonic on the neighborhood $\mathcal N$ of $\Gamma_{\mathrm{out}}$. In particular
\[
\tilde w\equiv 0\qquad\text{on }\mathcal N.
\]
\end{lemma}

\begin{proof}
We already know $\tilde w$ is $C^0$ on $\mathcal N$ and smooth and harmonic on
$\mathcal N\setminus\Gamma_{\mathrm{out}}$ (using Lemma~\ref{lem:w-harmonic} and
$\tilde w\equiv 0$ in $E_\infty\cap\mathcal N$).

For any test function $\varphi\in C_c^\infty(\mathcal N)$, we compute the distributional Laplacian.
Since $\tilde w$ equals $w$ on $\mathcal N\cap U$ and $0$ on $\mathcal N\cap E_\infty$:
\begin{align*}
\int_{\mathcal N}\tilde w\,\Delta\varphi\,dA
&=\int_{\mathcal N\cap U}w\,\Delta\varphi\,dA+\int_{\mathcal N\cap E_\infty}0\cdot\Delta\varphi\,dA\\
&=\int_{\mathcal N\cap U}w\,\Delta\varphi\,dA.
\end{align*}

By integration by parts in the region $\mathcal N\cap U$ (which is valid since $w$ is harmonic
and smooth away from $\Gamma_{\mathrm{out}}$):
\begin{align*}
\int_{\mathcal N\cap U}w\,\Delta\varphi\,dA
&=\int_{\mathcal N\cap U}\Delta w\cdot\varphi\,dA
+\int_{\Gamma_{\mathrm{out}}}\!\bigl(w\,\partial_\nu\varphi-\varphi\,\partial_\nu w\bigr)\,d\sigma\\
&=0+\int_{\Gamma_{\mathrm{out}}}\!\bigl(w\,\partial_\nu\varphi-\varphi\,\partial_\nu w\bigr)\,d\sigma,
\end{align*}
where $\nu$ is the outward unit normal to $U$ and we used that $w$ is harmonic
in $\mathcal N\cap U$ (so $\Delta w=0$ there).

The boundary integral vanishes:
\begin{itemize}
\item $w=0$ on $\Gamma_{\mathrm{out}}$ (Lemma~\ref{lem:cauchy-data}(b)), so the first term is $0$.
\item For the second term: since $\nabla u=0$ on $\Gamma_{\mathrm{out}}$ (Lemma~\ref{lem:cauchy-data}(a))
and $w=-y\partial_x u+x\partial_y u$ is a first-order linear combination of the components
of $\nabla u$ with coefficients depending only on $(x,y)$, the normal derivative
$\partial_\nu w$ vanishes on $\Gamma_{\mathrm{out}}$ (in the classical sense on the interior side,
and by continuity across the boundary).
\end{itemize}

Therefore $\int_{\mathcal N}\tilde w\,\Delta\varphi=0$ for every $\varphi\in C_c^\infty(\mathcal N)$, so
$\Delta\tilde w=0$ distributionally on $\mathcal N$. By Weyl's lemma, $\tilde w$ is a
classical harmonic (hence real-analytic) function on $\mathcal N$. Since $\tilde w$
vanishes on the nonempty open set $\mathcal N\cap E_\infty$, the identity principle yields
$\tilde w\equiv 0$ on $\mathcal N$.
\end{proof}

\subsection{From local vanishing to global rigidity}

\begin{lemma}\label{lem:w-zero-in-U}
Under (H1)--(H3), $w\equiv 0$ in every connected component $V$ of $U$ whose closure meets
$\Gamma_{\mathrm{out}}$. In particular, if $U$ is connected, $w\equiv 0$ in
$U^\circ\setminus\{0\}$.
\end{lemma}

\begin{proof}
Let $V$ be such a component. Lemma~\ref{lem:extension-harmonic} gives $w\equiv 0$ on the
open set $V\cap\mathcal N\ne\emptyset$. The set $V\setminus\{0\}$ is connected (a
connected planar open set remains connected after removing a point), $w$ is harmonic, hence
real-analytic, on $V\setminus\{0\}$, and vanishes on a nonempty open subset. By the identity
principle, $w\equiv 0$ on $V\setminus\{0\}$, hence on $V$ by continuity away from $0$.
\end{proof}

\begin{lemma}\label{lem:all-components}
Under (H1)--(H3), $w\equiv 0$ on all of $U\setminus\{0\}$.
\end{lemma}

\begin{proof}
By Lemma~\ref{lem:extension-harmonic}, $w\equiv 0$ on a neighborhood $\mathcal N$ of
$\Gamma_{\mathrm{out}}$ (the outer boundary of $U$).

Let $V$ be any connected component of $U$. We claim that $V\cap\mathcal N\ne\emptyset$.
Since $U$ is bounded and $V\subset U$, the closure $\overline V$ is compact. The outer
boundary $\Gamma_{\mathrm{out}}=\partial E_\infty$ separates $U$ from the unbounded exterior
$E_\infty$. Every bounded set in $\R^2$ has the property that each connected component
of the set must have a part of its boundary exposed to the unbounded exterior. Thus
$\overline V\cap\Gamma_{\mathrm{out}}\ne\emptyset$, and hence $V\cap\mathcal N\ne\emptyset$.

Now, the set $V\setminus\{0\}$ is a connected open set in $\R^2$ (a connected set remains
connected after removing a single point). By Lemma~\ref{lem:w-harmonic}, $w$ is harmonic
on $V\setminus\{0\}\subset\Omega=\R^2\setminus(\partial U\cup\{0\})$, so $w$ is real-analytic
there. Since $w$ vanishes identically on the nonempty open set $V\cap\mathcal N\subset
V\setminus\{0\}$, the identity principle for real-analytic functions gives $w\equiv 0$
on all of $V\setminus\{0\}$.

Since this holds for every connected component $V$ of $U$, and $U$ is the disjoint union
of its components, we conclude $w\equiv 0$ on $U\setminus\{0\}$.
\end{proof}

\begin{remark}
The proof of Lemma~\ref{lem:all-components} makes no appeal to nesting trees or induction
on connectivity; it uses only the analyticity of $w$ and the identity principle. This makes
the argument transparent for both connected and multiply-connected $U$.
\end{remark}

\subsection{Conclusion}

\begin{proof}[Proof of Theorem~\ref{thm:main}, $\Rightarrow$ direction]
Combining Lemmas \ref{lem:c-eq} and \ref{lem:all-components}, we have $c=\int_U f\,dA$ and
$\partial_\theta u\equiv 0$ on $U\setminus\{0\}$, so $u$ is radial on $U\setminus\{0\}$.
Also $u\equiv 0$ on $E_\infty$, which is radial (as complement of a relatively closed set,
once we know $U$ is radial; here we simply note $u|_{E_\infty}=0$ is radial trivially).
In any bounded component of $\R^2\setminus\overline U$ disjoint from the origin, $u$
is harmonic; as a harmonic function whose angular derivative
vanishes on the adjacent component of $U$ (hence on a full open arc of the boundary of
the bounded component) and which is $C^1$ across that boundary, $u$ is also radial there by
the same Cauchy uniqueness argument as above.

Thus $u$ is radial on $\R^2\setminus\{0\}$.

Returning to the distributional equation \eqref{eq:main}: both $\Delta u$ and $c\,\delta_0$
are rotationally invariant distributions (rotation invariance commutes with $\Delta$), so
\[
f(|x|)\1_U=\Delta u+c\,\delta_0
\]
is rotationally invariant as a distribution. Because $f(|x|)$ is a strictly positive
continuous function on $(0,\infty)$ and is itself rotationally invariant, the product
$f(|x|)\1_U$ being rotationally invariant forces $\1_U$ to be rotationally invariant as a
locally integrable function. Concretely, for any rotation $R_\theta$ about the origin, the
identity $f(|x|)\1_U(x)=f(|x|)\1_U(R_\theta x)$ holds almost everywhere in
$\R^2\setminus\{0\}$, and since $f(|x|)>0$ on that set, $\1_U(x)=\1_U(R_\theta x)$ almost
everywhere. Equivalently, the symmetric difference $U\triangle R_\theta U$ has Lebesgue
measure zero; since both sets are open, this symmetric difference is open with empty
interior, hence empty. Therefore $U=R_\theta U$ for every rotation, i.e., $U$ is
rotationally invariant about the origin.
\end{proof}

\section{The connected case}

\begin{corollary}\label{cor:connected}
Under the hypotheses of Theorem~\ref{thm:main}, if $U$ is connected then $U$ is either a
disk or an annulus centered at the origin.
\end{corollary}

\begin{proof}
A bounded connected open subset of $\R^2$ that is rotationally invariant about the origin
is determined by its intersection with a ray from the origin; this intersection must be an
open connected subset of $(0,\infty)$ (possibly with $0$ included, if $0\in U$), i.e., an
open interval. Rotating gives either an open disk $\{|x|<R\}$ (if $0\in U$) or an open
annulus $\{r_1<|x|<r_2\}$ (if $0\notin U$). Theorem~\ref{thm:main} gives the rotational
invariance; the shape classification follows.
\end{proof}

\section{Why this works where \texttt{main.tex}'s simply connected argument did not}

The approach in \texttt{main.tex} (Section~4) constructs a holomorphic function $W$ on $U$
whose boundary values are real, then invokes the maximum principle for $\Im W$ inside a
\emph{simply connected} $U$ to conclude $\Im W\equiv 0$. In multiply-connected $U$, the
maximum principle fails: the values of $\Im W$ on inner boundary components are not
controlled, making this approach unsuitable.

Our proof avoids this entirely by working with the real, first-order angular derivative
$w = L u = \partial_\theta u$ rather than with complex or holomorphic objects:

\begin{itemize}
\item[$(\alpha)$] \emph{Algebraic harmonicity:} $w$ is harmonic in $\Omega=\R^2\setminus
(\partial U\cup\{0\})$ via the commutation $\Delta w = L(\Delta u) = 0$, which holds for
purely algebraic reasons (rotation invariance of the right-hand side away from $\partial U$).
No topology is needed.

\item[$(\beta)$] \emph{Universal Cauchy data:} The Cauchy data of $w$ on $\Gamma_{\mathrm{out}}$
(any $C^1$ boundary, not necessarily a circle) is zero because:
$u\equiv 0$ in the unbounded exterior (Lemma~\ref{lem:vanish-Einfty}) $\Rightarrow$
$u=0,\nabla u=0$ on the boundary (Lemma~\ref{lem:cauchy-data}) $\Rightarrow$
$w=\partial_\theta u=0$ and $\partial_\nu w=0$ on the boundary.

\item[$(\gamma)$] \emph{Extension principle:} A harmonic function with zero Cauchy data
extends through a $C^1$ boundary as a harmonic function (Lemma~\ref{lem:extension-harmonic}).
This is a pure PDE fact requiring no topological assumptions beyond continuity.

\item[$(\delta)$] \emph{Propagation via analyticity:} Once $w=0$ near $\Gamma_{\mathrm{out}}$,
the fact that each connected component $V$ of $U$ has $V\cap\mathcal N\ne\emptyset$, combined
with the real-analyticity of $w$ on $V\setminus\{0\}$, forces $w\equiv 0$ throughout $V$
(Lemma~\ref{lem:all-components}). This uses only the analyticity of harmonic functions
and the identity principle, not simple connectivity.
\end{itemize}

\noindent\textbf{Key insight:} The conclusion ``$U$ is rotationally invariant''
(hence $\Gamma_{\mathrm{out}}$ is a circle) is a \emph{consequence} of the argument,
not a hypothesis. Nothing about $\Gamma_{\mathrm{out}}$ is assumed beforehand.

\begin{remark}
The entire proof uses only:
\begin{itemize}
\item Rotational symmetry of the right-hand side $f(|x|)\mathbf 1_U-c\,\delta_0$
\item $C^1$ regularity of $\partial U$ (for Green's formula)
\item Compact support of $u$ and continuity $u\in C^1(\R^2\setminus\{0\})$
\end{itemize}
No appeal to simple connectivity, multiply connectivity, moment identities, holomorphic
machinery, Bargmann/Hermite structures, or free-boundary theory is made. The theorem holds
for \emph{every} continuous positive radial weight $f:(0,\infty)\to\R$ with $f>0$.
\end{remark}

\end{document}
